\documentclass[12pt]{article}

\title{Simple Calculator}
\author{
	Evan M Drake \\
	\texttt{drakeev@butte.edu}
}


\usepackage{snow-preamble}

%%%%%%%%%%%%%%%%%%%%%%%%%%%%%%%%%%%%
\begin{document}
	\maketitle
	\tableofcontents

	\newpage

	%% body start
	\section{Interface}
		\indent The program will take two numbers, calculating some decimal value.
		That is any equation of the form $x + y$ should be evaluated,
			so long as $x$ \& $y$ are natural numbers.
		There will be no way for the user to inert values while the program is running.
		In particular the evaluated values must be in the code itself as two variables as below.

		\begin{lstlisting}[style=pseudocode, language=Python]
x = <left number>
y = <right number>
o = <operation>

def evaluate(x,y):
	return c
	# such that you use the correct
	# operation on x and y resulting in c
		\end{lstlisting}

	\section{Evaluation}
		\indent You should try to figure out a way to check what operation is desired.
		That is you may not have a different function for each binary arithmetic operation.
		You must find a way to dynamically detect what operation is desired.
		Hint: you might keep your third variable as a character and simply make a decision based on that character.
	
	\section{Testing}
		\indent you should be able to evaluate the expressions below:
		\begin{enumerate}
			\item $5 + 5$
			\item $1 + (-1)$
			\item $0 - (-1)$
			\item $\frac{0}{5}$
			\item $\frac{4}{2}$
			\item $2 * 1$
			\item $3 * 2$
		\end{enumerate}

	%% body finish

	\nocite{*}

	\newpage % bib
	\backmatter
	\addcontentsline{toc}{section}{References}
	\printbibliography

	\newpage % index
	\printindex

	\newpage % appendices
	\appendix

\end{document}

